\documentclass[11pt,a4paper]{article}
\usepackage[UTF8]{ctex}
\usepackage{amsmath,amssymb}
\usepackage{geometry}\geometry{margin=2cm}
\usepackage{graphicx,booktabs,array,longtable,multirow}
\usepackage{hyperref,xcolor,enumitem,caption,subcaption}
\hypersetup{colorlinks=true,linkcolor=blue,urlcolor=blue,citecolor=blue}
\setlength{\parskip}{0.3em}
\usepackage{fancyhdr,titlesec}
\pagestyle{fancy}\fancyhf{}\fancyhead[L]{\small 多语言内容审核模型技术报告}\fancyhead[R]{\small \thepage}
\title{\textbf{多语言内容审核模型\\技术报告与全面对比分析}}
\author{基于 microsoft/mdeberta-v3-base \\ 8 个国家 $\times$ 4 种训练方法 $\times$ 32+ 个模型}
\date{2026年6月12日}
\begin{document}
\maketitle
\tableofcontents
\newpage

% ════════════════════════════════════════════════════
\section{项目概述}

\subsection{项目背景}

随着社交媒体在非英语国家的快速普及，多语言内容审核成为全球化平台的核心挑战。
不同国家和地区有各自独特的违规定义：
沙特禁止 LGBTQ 内容，泰国严格保护王室声誉 (l\`{e}se-majest\'{e})，
南非有独特的种族隔离历史导致的严重种族主义 (k-word) 和排外主义问题，
新加坡有《煽动法令》和《维持宗教和谐法案》(MRHA) 监管种族/宗教不和谐言论。

本项目构建了一套覆盖 \textbf{8 个国家}的 AI 内容审核模型系统，
使用统一的多语言预训练基座 (mDeBERTa-v3-base)，
通过 4 种不同的微调策略适配各国的语言和法律要求。

\subsection{覆盖国家与语言}

\begin{table}[h]\centering
\caption{目标国家与语言}\label{tab:countries}
\begin{tabular}{@{}lllll@{}}\toprule
\textbf{代码} & \textbf{国家} & \textbf{语言} & \textbf{论坛/平台} & \textbf{数据量(合计)} \\\midrule
SA & 沙特阿拉伯 & 阿拉伯语 & Reddit r/saudiarabia & 85,824 \\
BR & 巴西 & 葡萄牙语 & Reddit r/brasil & 92,035 \\
MX & 墨西哥 & 西班牙语 & Reddit r/mexico & 64,013 \\
ID & 印尼 & 印尼语 & Reddit, Kompasiana & 40,675 \\
SG & 新加坡 & 英语/Singlish & HardwareZone, Reddit & 102,746 \\
TH & 泰国 & 泰语 & Pantip, Reddit & 20,768 \\
TR & 土耳其 & 土耳其语 & Reddit & 102,305 \\
ZA & 南非 & 英语/阿非利卡语等 & Reddit r/southafrica & 98,116 \\
\bottomrule
\end{tabular}
\end{table}

\subsection{任务定义}

\begin{enumerate}[itemsep=3pt]
\item \textbf{二分类 (Binary Classification)}:
输入一段社交媒体文本，输出该文本是否包含违规内容。
模型输出两个概率 $P(\text{safe})$ 和 $P(\text{violation})$，
取 $\arg\max$ 作为预测结果。
损失函数使用类别加权交叉熵以处理数据不均衡。

\item \textbf{多分类 (Multi-Class Classification)}:
输入一段社交媒体文本，输出该文本属于哪个具体的违规类别。
这是一个 8--9 类的单标签分类问题 (每条文本只属于一个类别)。
模型使用 softmax 输出每个类别的概率分布，
取最大概率类别作为预测结果，
并通过 per-class threshold calibration 优化各类别的精确率/召回率权衡。
\end{enumerate}

% ════════════════════════════════════════════════════
\section{模型选型}

\subsection{基座模型: mDeBERTa-v3-base}

我们选择 \textbf{microsoft/mdeberta-v3-base} 作为所有模型的统一基座:

\begin{enumerate}[itemsep=2pt]
\item \textbf{多语言覆盖}: mDeBERTa-v3 在 CC100 多语言语料上预训练，覆盖 100+ 语言，
包括本项目涉及的全部语言。
\item \textbf{解耦注意力机制 (Disentangled Attention)}:
DeBERTa 将内容嵌入和位置嵌入分离计算注意力权重，
在跨语言迁移任务上优于 XLM-RoBERTa 和 mBERT。
\item \textbf{适中的模型大小}: 279M 参数 (base 版本)，
全量微调需要 1.1GB/国，LoRA 适配器仅需约 600K 可训参数 (3MB)。
\item \textbf{MIT 许可证}: 可自由用于商业和研究。
\end{enumerate}

\subsection{四种微调方法}

\begin{enumerate}[itemsep=3pt]
\item \textbf{全量微调 (Full FT)}: 所有 279M 参数参与训练，精度最高，每国 1.1GB。
\item \textbf{LoRA (Low-Rank Adaptation)}:
在 attention 权重矩阵上添加低秩分解适配器 $W = W_0 + BA$，仅训练 0.21\% 参数 ($\sim$600K)。
LoRA 的核心思想是预训练权重的更新位于低秩子空间:
$h = W_0 x + \frac{\alpha}{r} \cdot BA x$，
其中 $B \in \mathbb{R}^{d \times r}$, $A \in \mathbb{R}^{r \times k}$,
$r \ll \min(d,k)$。适配器仅 $\sim$3MB。
\item \textbf{Multi-LoRA}: 一个共享基座 + 8 国独立的 LoRA 适配器 (r=8, $\sim$2MB/国)。
通过 \texttt{PeftModel.from\_pretrained} 热切换国家。
\item \textbf{Adapter}: 在 Transformer 层间插入 bottleneck adapter，
训练参数 0.5\%，适配器约 2.5MB。精度最低但训练最快。
\end{enumerate}

% ════════════════════════════════════════════════════
\section{数据集}

\subsection{二分类数据}

\input{_tables.tex}

注: Full FT 和 LoRA 使用 Nine\_country 数据; Adapter 和 Multi-LoRA 使用 data\_v2 数据。

% ════════════════════════════════════════════════════
\section{多分类标签体系}

\subsection{通用标签 (所有国家)}

以下 7 个标签在所有国家的多分类模型中共用，定义一致:

\begin{table}[h]\centering
\caption{通用违规标签定义}\label{tab:universal}
\begin{tabular}{@{}llp{6cm}@{}}\toprule
\textbf{标签} & \textbf{中文} & \textbf{定义} \\\midrule
Dangerous\_Content & 危险内容 & 恐怖主义、武器、自残/自杀、严重暴力、毒品交易、诈骗/钓鱼、性犯罪/CSAM、血腥内容 \\
Harassment & 骚扰 & 针对个人的侮辱、粗口、网络霸凌、人肉搜索、外貌羞辱。针对群体的攻击不属于此类 \\
Hate\_Speech & 仇恨言论 & 基于身份的攻击: 种族、宗教、性别、国籍。国家特有标签优先于通用 Hate\_Speech \\
Sexually\_Explicit\_Information & 色情信息 & 色情内容、性行为描述、约炮/性暗示。未成年相关内容归入 Dangerous\_Content \\
Politically\_Sensitive\_Topics & 政治敏感 & \textbf{非本国}政治内容: 外国选举、国际冲突、外国领导人 \\
Cybersecurity\_Malware & 网络安全 & 恶意软件、黑客攻击、钓鱼、DDoS。金融诈骗归入 Dangerous\_Content \\
safe & 安全 & 正常内容——这是最后的兜底标签，只有在确认不属于任何违规类别时才使用 \\
\bottomrule
\end{tabular}
\end{table}

\subsection{标签优先级规则}

当一条文本同时匹配多个标签时，按以下优先级选择:
\[
\text{Dangerous\_Content} > \text{国家特有标签} > \text{Hate\_Speech} > \text{Harassment} > \text{Sexually\_Explicit} > \text{Politically\_Sensitive} > \text{Cybersecurity\_Malware} > \text{safe}
\]

\subsection{各国特有标签}

每个国家在通用标签之外，有 2--3 个该国特有的违规标签。

\subsubsection{ID (印尼)}

\begin{table}[h]\centering
\caption{ID (印尼) 特有违规标签}
\begin{tabular}{@{}llp{6.5cm}@{}}\toprule
\textbf{标签} & \textbf{中文} & \textbf{详细定义与区分} \\\midrule
ID\_SARA & 印尼宗教种族 & 针对印尼族群(Suku)、宗教(Agama)、种族(Ras)、群体间(Antargolongan)的攻击。包括反华人(Cina pelit/Cina penguasa)、反爪哇人、反巴布亚人(Papua primitif)、政治身份攻击(cebong/kampret/kadrun)——这些政治身份词在印尼语境中属于群体攻击 \\
ID\_Blasphemy & 印尼亵渎/诽谤 & 1) 宗教亵渎 (Pasal 156a KUHP): 嘲讽安拉/先知/古兰经。2) UU ITE 电子诽谤 (Pasal 27-28): 社交媒体上的严重人格攻击、散布虚假指控毁坏名誉 \\
\bottomrule
\end{tabular}
\end{table}

\subsubsection{SG (新加坡)}

\begin{table}[h]\centering
\caption{SG (新加坡) 特有违规标签}
\begin{tabular}{@{}llp{6.5cm}@{}}\toprule
\textbf{标签} & \textbf{中文} & \textbf{详细定义与区分} \\\midrule
SG\_Racial\_Religious\_Harmony & 种族宗教和谐 & 新加坡煽动法令和 MRHA 下的种族/宗教不和谐。CMIO 种族攻击。反移民种族言论(FT/Bangla/Pinoy/PRC/Ah Tiong) \\
SG\_Vulgarity\_Singlish & 粗俗 Singlish & 极端粗俗的 Singlish 侮辱: chee bai/cibai, kan ni na/kNN, lan jiao。关键区分: 日常 Singlish 语气词(lah, lor, leh)绝不算违规 \\
\bottomrule
\end{tabular}
\end{table}

\subsubsection{TH (泰国)}

\begin{table}[h]\centering
\caption{TH (泰国) 特有违规标签}
\begin{tabular}{@{}llp{6.5cm}@{}}\toprule
\textbf{标签} & \textbf{中文} & \textbf{详细定义与区分} \\\midrule
TH\_Lese\_Majeste & 冒犯君主 & 刑法第 112 条。对王室成员的侮辱、嘲讽、威胁。泰国最严重的网络言论罪行，即使是隐晦的批评也可能被归入此类 \\
TH\_Political\_Instigation & 政治煽动 & 泰国国内政治煽动: 推翻政府号召、军事政变呼吁、选举舞弊指控、暴力抗议号召。非泰国政治归入 Politically\_Sensitive\_Topics \\
\bottomrule
\end{tabular}
\end{table}

\subsubsection{ZA (南非)}

\begin{table}[h]\centering
\caption{ZA (南非) 特有违规标签}
\begin{tabular}{@{}llp{6.5cm}@{}}\toprule
\textbf{标签} & \textbf{中文} & \textbf{详细定义与区分} \\\midrule
ZA\_Severe\_Racism & 严重种族主义 & K-word (kaffir)——南非最严重的种族侮辱词，PEPUDA 法案下的刑事犯罪。还包括: 种族隔离侮辱词(hotnot, coolie)、种族隔离怀旧、种族优越论。一旦出现 k-word，不论上下文一律归入此类 \\
ZA\_Xenophobia & 排外主义 & 针对南非境内非洲外国人的排外攻击。核心特征: makwerekwere(贬损词)、要求驱逐津巴布韦人/尼日利亚人、Operation Dudula 言论。南非在 2008/2015/2019/2022 年经历了多轮致命排外暴力 \\
\bottomrule
\end{tabular}
\end{table}

\subsubsection{SA (沙特)}

\begin{table}[h]\centering
\caption{SA (沙特) 特有违规标签}
\begin{tabular}{@{}llp{6.5cm}@{}}\toprule
\textbf{标签} & \textbf{中文} & \textbf{详细定义与区分} \\\midrule
SA\_LGBTQ\_Content & LGBTQ 内容 & 沙特法律禁止的 LGBTQ 相关内容。在沙特语境下即使是中性的 LGBTQ 讨论也可能归入此类 \\
SA\_Religious\_Violation & 宗教违规 & 宗教亵渎、叛教(ridda)、对伊斯兰教法的侮辱、嘲讽宗教仪式 \\
SA\_State\_Security\_Royalty & 国家安全/王室 & 反政府言论、反王室家族(Al Saud)言论、号召抗议/示威 \\
\bottomrule
\end{tabular}
\end{table}

\subsubsection{MX (墨西哥)}

\begin{table}[h]\centering
\caption{MX (墨西哥) 特有违规标签}
\begin{tabular}{@{}llp{6.5cm}@{}}\toprule
\textbf{标签} & \textbf{中文} & \textbf{详细定义与区分} \\\midrule
MX\_Narco\_Culture & 毒品文化 & 赞美/美化毒品卡特尔(cartel)、毒枭(Chapo/Mencho)、毒品民谣(narcocorridos)。墨西哥已有 40 万+ 卡特尔相关死亡 \\
MX\_Gender\_Violence & 性别暴力 & 杀害女性(femicide)赞美/正当化、受害者归咎(se lo merec\'ia/她活该)、大男子主义暴力威胁 \\
\bottomrule
\end{tabular}
\end{table}

\subsubsection{BR (巴西)}

\begin{table}[h]\centering
\caption{BR (巴西) 特有违规标签}
\begin{tabular}{@{}llp{6.5cm}@{}}\toprule
\textbf{标签} & \textbf{中文} & \textbf{详细定义与区分} \\\midrule
BR\_Structural\_Racism & 结构性种族主义 & 严重反黑人种族主义。macaco(猴子)在巴西语境中是对黑人的严重侮辱。巴西有 350+ 年奴隶制历史(1888 年最后废除)——即使是"玩笑式"种族主义(racismo recreativo)也属严重违规 \\
BR\_Political\_Extremism & 政治极端主义 & 反民主极端主义: interven\c{c}\~{a}o militar j\'{a}(立即军事干预)、2022 选举否认、攻击 STF/国会/TSE。2023 年 1 月 8 日后被视为对民主的直接威胁 \\
\bottomrule
\end{tabular}
\end{table}

\subsubsection{TR (土耳其)}

\begin{table}[h]\centering
\caption{TR (土耳其) 特有违规标签}
\begin{tabular}{@{}llp{6.5cm}@{}}\toprule
\textbf{标签} & \textbf{中文} & \textbf{详细定义与区分} \\\midrule
TR\_Political\_Extremism & 政治极端主义 & 反政府极端言论、政变号召、恐怖组织(如 PKK)赞美、煽动分裂主义 \\
TR\_Religious\_Extremism & 宗教极端主义 & 宗教极端言论、教派攻击(Alevi/Sunni)、利用宗教煽动暴力 \\
\bottomrule
\end{tabular}
\end{table}

% ════════════════════════════════════════════════════
\section{各国家多分类训练数据分布}

多分类训练数据由三部分组成:
\textbf{data\_v2 标签映射}(基础)、\textbf{Nine\_country safe 样本}(补充安全内容)、
\textbf{NIM API 生成}(补充弱势类别)。

% Include the label distribution tables
\input{_tables.tex}

% ════════════════════════════════════════════════════
\section{训练方法详述}

\subsection{全量微调 (Full Fine-Tuning)}

全量微调是精度最高的方法——所有 279M 参数参与训练。

\textbf{训练配置:}
\begin{verbatim}
learning_rate = 1e-5
epochs = 5 (early stopping patience=3, monitor=eval_loss)
per_device_train_batch_size = 4
gradient_accumulation_steps = 4  # effective batch = 16
max_grad_norm = 0.5
weight_decay = 0.01, warmup_ratio = 0.1
lr_scheduler = "cosine", max_length = 256
precision = FP32
loss = CrossEntropyLoss(weight=class_weights)
\end{verbatim}

由于二分类数据中 safe/violation 比例不均衡，使用类别加权的交叉熵损失:
\[ \mathcal{L} = -\frac{1}{N}\sum_{i=1}^N w_{y_i} \log P(y_i|x_i) \]
其中 $w_0 = N/(2N_0)$ (safe 权重), $w_1 = N/(2N_1)$ (violation 权重)。

\subsection{LoRA 微调 (Low-Rank Adaptation)}

LoRA 通过在预训练权重上添加低秩分解矩阵实现参数高效微调:
\[ h = W_0 x + \frac{\alpha}{r} \cdot BA x \]
其中 $B \in \mathbb{R}^{d \times r}$, $A \in \mathbb{R}^{r \times k}$, rank $r=16 \ll \min(d,k)$。

\textbf{训练配置 (二分类):}
\begin{verbatim}
LoraConfig(r=16, lora_alpha=32, lora_dropout=0.1,
           target_modules=["query_proj", "value_proj"])
learning_rate = 5e-5, epochs = 5
per_device_train_batch_size = 16
precision = BF16
\end{verbatim}

\textbf{训练配置 (多分类):}
\begin{verbatim}
LoraConfig(r=16, lora_alpha=32, lora_dropout=0.1)
learning_rate = 2e-4, epochs = 3
per_device_train_batch_size = 16
max_length = 256
label_smoothing = 0.05
\end{verbatim}

\subsection{Label Smoothing 损失 (多分类)}

Label Smoothing 将 one-hot 标签"软化"为平滑分布，防止过拟合:
\[ \mathcal{L}_{LS} = -\sum_{c=1}^C \tilde{y}_c \log P(c|x) \]
其中 $\tilde{y}_c = (1-\epsilon) \cdot y_c + \epsilon/C$, $\epsilon = 0.05$。

Label Smoothing 在几乎所有多分类模型上稳定优于普通交叉熵 (+0.01--0.02 Macro F1)。

\subsection{Threshold Calibration (后处理优化)}

训练完成后，为每个类别独立优化分类阈值以最大化 Macro F1:
\[ \hat{y} = \arg\max_c \frac{P(c|x)}{t_c} \]

优化过程:
\begin{enumerate}
\item 对每个类别，在 $[0.02, 0.98]$ 以 0.02 步长搜索最优 $t_c$ (最大化 Macro F1)
\item 使用 Nelder-Mead 算法联合优化所有阈值
\end{enumerate}

Threshold calibration 在所有模型上稳定提供 +0.01--0.02 的提升。

\subsection{LSD 蒸馏 (已弃用)}

ID v2 中尝试了 Label Smoothing Distillation:
\[ \mathcal{L} = \alpha \cdot \text{CE}(ls=0.05) + (1-\alpha) \cdot \text{KL}(P_T/T \| P_S/T) \cdot T^2 \cdot w \]
($T=3.0$, $\alpha=0.4$)。提升有限 (+0.005)，训练慢 5 倍 (KL 散度开销)，不推荐。

% ════════════════════════════════════════════════════
\section{二分类模型结果}

\subsection{8 国 $\times$ 4 种方法对比}

图~\ref{fig:bin-all} 展示了 8 个国家 $\times$ 4 种方法的全面对比。

\begin{figure}[h]\centering
\includegraphics[width=\textwidth]{01_binary_comparison_cn.png}
\caption{8 国 $\times$ 4 种方法二分类 Macro F1 对比。全量微调 (绿) 在所有国家都表现最优，平均领先 LoRA +0.044}\label{fig:bin-all}
\end{figure}

\subsection{全量微调 (Full FT) -- 8 国完整结果}

% Table will be inserted from _tables.tex

\subsection{LoRA 微调 -- 8 国完整结果}

% Table will be inserted from _tables.tex

\subsection{四种方法对比}

\begin{table}[h]\centering
\caption{四种方法 8 国平均值对比}\label{tab:method-compare}
\begin{tabular}{@{}lrrrl@{}}\toprule
\textbf{方法} & \textbf{Avg Acc} & \textbf{Avg Macro F1} & \textbf{每国大小} & \textbf{适用场景} \\\midrule
Full FT & 0.853 & \textbf{0.853} & 1.1GB & 生产环境，精度优先 \\
LoRA (r=16) & 0.814 & 0.809 & 3MB & 快速部署，接近 FT \\
Multi-LoRA (r=8) & 0.860 & 0.796 & 2MB & 多国服务，省空间 \\
Adapter & 0.849 & 0.769 & 2.5MB & 快速原型 \\
\bottomrule
\end{tabular}
\end{table}

\subsection{关键发现}

\begin{enumerate}[itemsep=2pt]
\item 全量微调在所有 8 国中都是最优方法，平均领先 LoRA +0.044 Macro F1。
\item 新加坡 (SG) 表现最佳 (Full FT Acc 0.919, 数据量 102K)。
\item Multi-LoRA 违规召回率最高 (SA: 0.983)，但 Clean Recall 低——倾向将内容判为违规。
\item 墨西哥 (MX) 的 LoRA 退化最严重 (-0.079)，可能因西班牙语数据分布与预训练语料差异大。
\end{enumerate}

% ════════════════════════════════════════════════════
\section{多分类模型结果}

\subsection{整体排名}

图~\ref{fig:mc-rank} 和图~\ref{fig:heatmap} 展示了多分类模型对比。

\begin{figure}[h]\centering
\includegraphics[width=0.85\textwidth]{02_multiclass_ranking_cn.png}
\caption{多分类模型 Macro F1 排名。泰国 (0.744) 高于印尼 (0.732)，可能因为标签更少 (8 vs 9) 且数据噪音更低}\label{fig:mc-rank}
\end{figure}

\begin{figure}[h]\centering
\includegraphics[width=\textwidth]{03_multiclass_heatmap_cn.png}
\caption{各标签 $\times$ 各国 F1 热力图。绿色=高 F1，红色=低 F1。NIM 生成的标签得分普遍较高；国家特有标签得分较低}\label{fig:heatmap}
\end{figure}

\begin{table}[h]\centering
\caption{多分类模型结果汇总}\label{tab:mc-summary}
\begin{tabular}{@{}lrrrl@{}}\toprule
\textbf{国家} & \textbf{标签数} & \textbf{训练量} & \textbf{Macro F1} & \textbf{训练方法} \\\midrule
沙特 (SA) & 8 & $\sim$15K & \textbf{0.829} & 全量微调 \\
泰国 (TH) & 8 & 17,751 & \textbf{0.744} & LoRA LS + cal \\
印尼 (ID v6) & 9 & 45,753 & \textbf{0.732} & LoRA LS + cal \\
新加坡 (SG) & 9 & 25,543 & \textbf{0.704} & LoRA LS + cal \\
南非 (ZA v1) & 9 & 21,329 & \textbf{0.589} & LoRA LS + cal \\
\bottomrule
\end{tabular}
\end{table}

\subsection{各国各类别详细结果}

% Per-class tables from _tables.tex

% ════════════════════════════════════════════════════
\section{印尼 (ID) 模型演进}

印尼模型经历了最完整的 6 版本迭代优化。

\begin{figure}[h]\centering
\includegraphics[width=0.9\textwidth]{04_id_evolution_cn.png}
\caption{ID 模型从 v1 到 v6 的演进路线。每次 NIM 数据生成都带来显著提升}\label{fig:id-evo}
\end{figure}

\textbf{演进关键节点:}
\begin{enumerate}[itemsep=2pt]
\item \textbf{v1 (38K, 8-class)}: Full FT (0.692) 和 LoRA (0.690) 基线。
\item \textbf{v2 (38K, 9-class)}: LSD 蒸馏，0.684。训练慢 5 倍，弃用。
\item \textbf{v3}: 3 个 NIM 陪审员模型验证 300 个混淆样本，修正 ID\_SARA/Harassment 标签。
\item \textbf{v4}: NIM 生成 2,500 条 ID\_SARA。ID\_SARA F1: 0.42 $\to$ 0.62 (+0.20)。
\item \textbf{v5 (42K)}: 2-seed LoRA LS ensemble + threshold cal，F1 0.709。
\item \textbf{v6 (46K)}: +2,000 Dangerous +1,500 Hate\_Speech，\textbf{F1 0.732}。
\end{enumerate}

\begin{figure}[h]\centering
\includegraphics[width=0.85\textwidth]{06_data_augmentation_impact_cn.png}
\caption{NIM 数据生成对弱势类别的提升效果}\label{fig:augment}
\end{figure}

% ════════════════════════════════════════════════════
\section{综合仪表盘}

\begin{figure}[h]\centering
\includegraphics[width=\textwidth]{05_dashboard_cn.png}
\caption{综合仪表盘——6 个维度全面对比}\label{fig:dashboard}
\end{figure}

% ════════════════════════════════════════════════════
\section{关键发现与经验总结}

\begin{enumerate}[itemsep=3pt]
\item \textbf{全量微调是二分类最佳方法}: 在所有 8 国中 Full FT 都显著优于 PEFT 方法。
\item \textbf{NIM 数据生成是最有效的改进手段}: ID\_SARA +0.23, Dangerous +0.15, Hate\_Speech +0.09。
\item \textbf{Label Smoothing (0.05) 稳定优于普通 CE}: +0.01--0.02 的免费提升。
\item \textbf{Threshold Calibration 稳定提供 +0.01--0.02}: 对所有模型有效。
\item \textbf{LSD 蒸馏不推荐}: 训练慢 5 倍，提升不足 0.01。
\item \textbf{生成数据可能过拟合}: ZA Political (0.99) 和 TH Hate\_Speech (0.96) 的接近完美 F1
暗示模型学会了 NIM 生成的模式而非真正违规特征。
\item \textbf{LoRA 适配器降低 460 倍部署成本}: 2.4MB vs 1.1GB。
\end{enumerate}

% ════════════════════════════════════════════════════
\section{未来工作}

\begin{enumerate}[itemsep=2pt]
\item \textbf{ZA 重训 (v2)}: 补充 ZA\_Xenophobia (+2,500) 和 ZA\_Severe\_Racism (+2,500) 后重训。
\item \textbf{BR/MX/TR 多分类}: data\_v2 标签映射已就绪，可快速启动。
\item \textbf{Multi-LoRA 多分类升级}: 一个基座服务 8 国多类违规检测。
\item \textbf{消除 NIM 生成数据的分布偏差}: 用陪审团验证过滤低质量样本。
\end{enumerate}

\end{document}
